\cotefile{157-2}
\scanpage{428}
\mspage{167}
% archive sheet number: 427

c'est exactement le même argument que tantôt, \deleted{\unsure{mais}} je présume. Il y a cependant une complication technique. Nous \deleted{\illegible{1}} avions pris, pour définir

\[
\uline{t}\colon\mathbb{N}^{\unsure{\circ|o}}\longrightarrow\uline{\unsure{\mathrm{Ty}|\mathrm{T}y}}\qquad \uline{t}=(t_n),
\]

le choix

\[
t_n=(\Delta_+\Delta_-)^{2n},
\]

\unsure{où|ou} si on préfère

\[
t_n=(\Delta_+\Delta_{\unsure{+|-}})^n,
\]

car tous \deleted{\unsure{les}} les $t_n$ sont \unsure{à|:} \unsure{direction|de variation} initiale $+$, terminale $-$, et ce sont des types réduits. Si on prend $\Delta_+t_n\Delta_-$, on trouve un type non réduit, qui n'est pas parmi les $t_n$. Il ne servirait à rien de changer la définition des $t_n$ en $t_n=(\Delta_-\Delta_+)^{\text{\unsure{qq chose}}}$, car cette fois $\Delta_+t_n\Delta_-$ sera bien réduit, mais n'aura pas les bons signes initial et terminal. Donc pour définir $\widetilde c$ en \unsure{termes|terme} de $c$, comme objet de $\unsure{\mathbb{C}|\mathcal{C}}_{\mathbb{N}^{\circ}}$, \deleted{\unsure{il}} il faudra, si $c$ est de type $\tau$\inserted{$=t_n$}, prendre $\widetilde c$ de type $\Delta_+\Delta_-\tau\Delta_+\Delta_-=t_{n+1}$ (en prenant $t_{n+1}=(\Delta_+\Delta_-)^{\unsure{2n|2m}})$, et de \unsure{$=$|même} pour définir $\widetilde{c'}$ en terme de $c'$. Les \unsure{dispositions|définitions} pertinentes \deleted{\unsure{sont décrites}} le passage de $(c',w')$ à $(\widetilde{c'},\widetilde{\unsure{w'|w_*}})$ est décrit dans la figure
